%! Trigonometric Equations and Inequalities — Unit 3, Lesson 3.10 — Student Packet Cover Sheet
\documentclass[10pt]{article}
\usepackage{precalculus-article}
\usepackage{precalculus-boxes}
\usepackage{ltablex}
\keepXColumns

\begin{document}

% Full-bleed navy banner
\begin{tikzpicture}[remember picture, overlay]
  \fill[navy] (current page.north west)
    rectangle ([yshift=-0.9in]current page.north east);
\end{tikzpicture}
\vspace{-0.8in}

\begin{center}
  {\color{white}\textbf{\LARGE Precalculus}}\\[4pt]
  {\color{skymid}\large\textbf{Unit 3: Trigonometric and Polar Functions}}\\[6pt]
  {\color{white}\textbf{\large Lesson 3.10 \quad Trigonometric Equations and Inequalities}}
\end{center}

\vspace{0.22in}
\namedateperiod
\vspace{0.18in}

\begin{learningtargetbox}
\small
\begin{enumerate}[leftmargin=1.5em, itemsep=5pt]
  \item I can solve a trigonometric equation by isolating the trig function,
    applying an inverse trig function to find the principal value, and using
    periodicity to write the general solution. \hfill\textit{(LO 3.10.A)}
  \item I can explain why trigonometric equations have infinitely many solutions
    when the domain is unrestricted, and identify both solution families.
    \hfill\textit{(LO 3.10.A, EK 3.10.A.2)}
  \item I can apply a domain restriction implied by a contextual scenario to
    select only the solutions that are valid in that context.
    \hfill\textit{(LO 3.10.A, EK 3.10.A.3)}
\end{enumerate}
\end{learningtargetbox}

\vspace{0.14in}

\begin{tocbox}
\small
\renewcommand{\arraystretch}{1.5}
\begin{tabularx}{\linewidth}{@{} c l X r @{}}
\rowcolor{sky}
\textbf{\#} & \textbf{Component} & \textbf{Description} & \textbf{Score} \\
\midrule
1 & Warm-Up        & Spiral review: inverse trig values, unit circle, isolating trig functions & \blank{1.2cm} \\
2 & Guided Notes   & General solutions, inverse trig procedure, contextual domain, inequalities & \blank{1.2cm} \\
3 & Group Activity & Solve trig equations and inequalities (tiered R/A/E) & \blank{1.2cm} \\
4 & Exit Ticket    & Solve $\sin\theta = -\tfrac{\sqrt{3}}{2}$; explain infinitely many solutions & \blank{1.2cm} \\
5 & Homework       & Mixed: general solutions, context, AP MC, quadratic substitution extension & \blank{1.2cm} \\
\midrule
  & \multicolumn{2}{r}{\textbf{Total}} & \blank{1.2cm} \\
\end{tabularx}
\end{tocbox}

\end{document}
